\documentclass[11pt]{article}
\usepackage[margin=1in]{geometry}
\usepackage{amsmath,amssymb,amsthm,booktabs,array,longtable,enumitem,hyperref,microtype}
\usepackage[T1]{fontenc}
\hypersetup{colorlinks=true,linkcolor=blue,citecolor=blue,urlcolor=blue}

% Identification strip in the left margin of every page (metadata only).
% The DOI identifies the published Zenodo record; scientific content is unchanged.
\usepackage{eso-pic,graphicx,xcolor}
\definecolor{stripdoi}{RGB}{62,151,208}
\definecolor{striplinks}{RGB}{85,158,110}
\newcommand{\ResearchMarginStrip}{%
  \begingroup\sffamily\fontsize{7}{9}\selectfont
  \href{https://doi.org/10.5281/zenodo.23533046}{\textcolor{stripdoi}{10.5281/zenodo.23533046}}%
  \hspace{1.6em}\textcolor{black}{v1 [NT] 21 Oct 2026}%
  \hspace{1.6em}%
  \href{https://orcid.org/0009-0004-7023-0298}{\textcolor{striplinks}{ORCID}}%
  \textcolor{black}{ / }%
  \href{https://github.com/FarhadBanazadeh}{\textcolor{striplinks}{GITHUB}}%
  \endgroup
}
\AddToShipoutPictureFG{%
  \AtPageLowerLeft{%
    \put(15,\LenToUnit{.5\paperheight}){%
      \makebox(0,0){\rotatebox{90}{\ResearchMarginStrip}}%
    }%
  }%
}

\newtheorem{theorem}{Theorem}[section]
\newtheorem{proposition}[theorem]{Proposition}
\newtheorem{lemma}[theorem]{Lemma}
\newtheorem{corollary}[theorem]{Corollary}
\newtheorem{conjecture}[theorem]{Conjecture}
\theoremstyle{remark}\newtheorem{remark}[theorem]{Remark}\newtheorem{example}[theorem]{Example}
\newcommand{\vfive}{\nu_5}
\newcommand{\N}{\mathbb N}
\title{A $6^r$-Scaled Family of the Reduced $\lfloor 6k/5\rfloor+1$ Collatz-Type Map\\[3pt]
\large Exact Algebraic Conjugacy, Floor-Form Representation, $5$-Adic Reduction Invariance, Cycle and Inverse-Tree Transport, and Transfer of the $10^{11}$ Finite Verification}
\author{Farhad Banazadeh\\\texttt{fn.bana@hotmail.com}\\ORCID: 0009-0004-7023-0298}
\date{21 October 2026}
\begin{document}\maketitle
\begin{abstract}
We develop the $6^r$-scaled family associated with the reduced map generated by
$G_0(n)=\lfloor 6n/5\rfloor+1$ and complete removal of factors of $5$.
Writing $c(n)=5-(n\bmod 5)$ with the convention $c(n)=5$ when $5\mid n$, one has
$G_0(n)=(6n+c(n))/5$. For every integer $r\ge0$ we introduce the scaled lattice
$D_r=6^rD$, where $D=\{n\ge1:5\nmid n,\ 6\nmid n\}$ is the closed post-transient domain of the base system, and define
\[
 H_r(x)=\frac{6x+6^r c(x)}5,\qquad T_r(x)=\frac{H_r(x)}{5^{\vfive(H_r(x))}}.
\]
Because $6^r\equiv1\pmod5$ and $5\nmid6^r$, multiplication by $6^r$ preserves the residue branch and contributes no $5$-adic valuation. We prove the exact conjugacy
\[
T_r(6^r n)=6^rT_0(n),\qquad T_r^j(6^r n)=6^rT_0^j(n).
\]
The generalized initial quotient also has the requested floor representation
\[
H_r(x)=\left\lfloor\frac{6x}{5}\right\rfloor+1+
\frac{6^r-1}{5}\Bigl[5\bigl(\lfloor6x/5\rfloor+1\bigr)-6x\Bigr]
\]
for $5\nmid x$, with the bracket selecting the appropriate coefficient in $\{1,2,3,4\}$.
Consequently the complete directed functional graph on the scaled domain, the period-four attractor, cycle periods, predecessor trees, valuation words, first-descent times, elementary-step counts, peak positions, and normalized peak ratios are transported exactly. Absolute states and peaks are multiplied by $6^r$.
The completed base first-descent verification through $10^{11}$ therefore transfers, without a new $10^{11}$-start scan, to the corresponding scaled set at every fixed $r$. This is an exact theorem on corresponding scaled states, not an independent exhaustive computation at each scale. The global convergence question remains equivalent to the unresolved base problem.
\end{abstract}

\section{Introduction}
The base reduced floor map is defined by
\[
G_0(n)=\left\lfloor\frac{6n}{5}\right\rfloor+1,
\qquad T_0(n)=\frac{G_0(n)}{5^{\vfive(G_0(n))}}.
\]
A completed finite computation for the base system exhaustively treated every start $1\le n\le10^{11}$ and established strict first descent for every $n\ge2$ in that interval, hence entry into the exact four-cycle $(1,2,3,4)$ by strong induction. The purpose of the present paper is structural: to determine the precise dynamical meaning of replacing the residue correction $1,2,3,4$ by $6^r,2\cdot6^r,3\cdot6^r,4\cdot6^r$.

The key point is that the generalized affine expression becomes a true scaled copy of the base dynamics only on the corresponding scaled lattice. On that domain, the factor $6^r$ can be pulled through the numerator and survives complete $5$-adic reduction. This yields an exact conjugacy rather than a heuristic similarity.

We distinguish throughout: (i) exact algebraic statements valid for all $r\ge0$; (ii) finite results inherited from the completed base computation; and (iii) probabilistic interpretations inherited only after normalization. No fresh $10^{11}$-start scan is claimed for each scale.

\section{The base reduced floor map}
Write $n=5q+s$, $0\le s\le4$, and define
\[
c(n)=\begin{cases}5,&s=0,\\5-s,&s=1,2,3,4.\end{cases}
\]
Then
\begin{equation}\label{eq:base-affine}
G_0(n)=\frac{6n+c(n)}5=\left\lfloor\frac{6n}{5}\right\rfloor+1.
\end{equation}
Let $e_0(n)=\vfive(G_0(n))$ and
\[
T_0(n)=G_0(n)/5^{e_0(n)}.
\]
The base study proves that after one iteration every orbit lies in
\[
D=\{m\ge1:5\nmid m,\ 6\nmid m\},
\]
and $T_0(D)\subseteq D$. On $D$, $c(n)\in\{1,2,3,4\}$. The exact cycle
\[
C=(1,2,3,4),\qquad1\mapsto2\mapsto3\mapsto4\mapsto1,
\]
is contained in $D$.

The completed base certificate covers $10^{11}$ starting values. Its headline records used later are shown in Table~\ref{tab:base}.
\begin{table}[ht]\centering
\caption{Base finite-verification record inherited from the completed base computation.}\label{tab:base}
\begin{tabular}{lr}\toprule
Quantity & Base result\\\midrule
Starts tested & $100,000,000,000$\\
Maximum first-descent time & $422$\\
First start attaining it & $55,335,461,598$\\
Maximum elementary count & $470$\\
Maximum observed extra valuation & $16$\\
Largest pre-descent state & $8,669,800,813,122,806,007,604$\\
Start producing largest peak & $37,266,730,278$\\
Mean first-descent time & $6.18073627485$\\\bottomrule
\end{tabular}
\end{table}

\section{Definition of the $6^r$-scaled family}
Fix $r\ge0$ and set
\[
D_r=6^rD=\{6^rn:n\in D\}.
\]
For $x\in D_r$ define
\begin{equation}\label{eq:Hr}
H_r(x)=\frac{6x+6^r c(x)}5,
\end{equation}
and
\begin{equation}\label{eq:Tr}
e_r(x)=\vfive(H_r(x)),\qquad T_r(x)=\frac{H_r(x)}{5^{e_r(x)}}.
\end{equation}
Since $x\in D_r$ implies $5\nmid x$, the correction $c(x)$ is one of $1,2,3,4$. Thus \eqref{eq:Hr} is precisely one of
\[
\frac{6x+6^r}{5},\quad \frac{6x+2\cdot6^r}{5},\quad
\frac{6x+3\cdot6^r}{5},\quad \frac{6x+4\cdot6^r}{5},
\]
with the coefficient forced by $x\bmod5$.

\section{Exact floor-function representation}
The residue selector can itself be written using the base floor expression.
\begin{lemma}[Floor residue selector]\label{lem:selector}
For $5\nmid x$,
\[
c(x)=5\left(\left\lfloor\frac{6x}{5}\right\rfloor+1\right)-6x.
\]
\end{lemma}
\begin{proof}
Write $x=5q+s$ with $s\in\{1,2,3,4\}$. Then $\lfloor6x/5\rfloor=6q+s$, so the right side is $5(6q+s+1)-6(5q+s)=5-s=c(x)$.
\end{proof}

\begin{theorem}[Generalized floor form]\label{thm:floor}
For $x\in D_r$,
\begin{align}
H_r(x)&=\left\lfloor\frac{6x}{5}\right\rfloor+1+
\frac{6^r-1}{5}\,c(x)\label{eq:floor-short}\\
&=\left\lfloor\frac{6x}{5}\right\rfloor+1+
\frac{6^r-1}{5}\left[5\left(\left\lfloor\frac{6x}{5}\right\rfloor+1\right)-6x\right].\label{eq:floor-full}
\end{align}
\end{theorem}
\begin{proof}
From \eqref{eq:base-affine}, $(6x+c(x))/5=\lfloor6x/5\rfloor+1$. Therefore
\[
\frac{6x+6^rc(x)}5=\frac{6x+c(x)}5+\frac{(6^r-1)c(x)}5.
\]
Since $6^r\equiv1\pmod5$, the second coefficient is integral. Lemma~\ref{lem:selector} gives \eqref{eq:floor-full}.
\end{proof}
At $r=0$ the correction term vanishes and the base floor map is recovered exactly.

\section{Scaling lemmas}
\begin{lemma}[Residue preservation]\label{lem:res}
For every $n\in D$ and $r\ge0$,
\[
6^rn\equiv n\pmod5,\qquad c(6^rn)=c(n).
\]
\end{lemma}
\begin{proof}Since $6\equiv1\pmod5$, also $6^r\equiv1\pmod5$.\end{proof}

\begin{lemma}[Exact affine factorization]\label{lem:factor}
For every $n\in D$,
\[
H_r(6^rn)=6^rG_0(n).
\]
\end{lemma}
\begin{proof}
By Lemma~\ref{lem:res},
\[
H_r(6^rn)=\frac{6(6^rn)+6^rc(n)}5=6^r\frac{6n+c(n)}5.
\]
\end{proof}

\begin{lemma}[$5$-adic valuation invariance]\label{lem:v5}
For every nonzero integer $m$,
\[
\vfive(6^rm)=\vfive(m).
\]
Consequently $e_r(6^rn)=e_0(n)$.
\end{lemma}
\begin{proof}The scale factor $6^r$ is coprime to $5$.\end{proof}

\begin{corollary}[Scaled closure]
$T_r(D_r)\subseteq D_r$.
\end{corollary}

\section{Exact conjugacy and functional-graph isomorphism}
Define $\Phi_r:D\to D_r$ by $\Phi_r(n)=6^rn$.
\begin{theorem}[One-step conjugacy]\label{thm:conj}
For every $r\ge0$ and $n\in D$,
\[
T_r(6^rn)=6^rT_0(n),\qquad T_r\circ\Phi_r=\Phi_r\circ T_0.
\]
\end{theorem}
\begin{proof}Combine Lemmas~\ref{lem:factor} and \ref{lem:v5}.\end{proof}

\begin{theorem}[All iterates]
For every $j\ge0$,
\[
T_r^j(6^rn)=6^rT_0^j(n).
\]
\end{theorem}
\begin{proof}Induction on $j$ using Theorem~\ref{thm:conj}.\end{proof}

\begin{corollary}[Mutual conjugacy of scale levels]
For $r,s\ge0$, the map $\Psi_{r,s}(6^rn)=6^sn$ is a bijection $D_r\to D_s$ satisfying
$T_s\circ\Psi_{r,s}=\Psi_{r,s}\circ T_r$.
\end{corollary}
Thus the directed functional graphs are isomorphic after normalization.

\section{The scaled four-cycle and cycle transport}
\begin{corollary}[Scaled attractor cycle]
The base cycle transports to
\[
C_r=(6^r,2\cdot6^r,3\cdot6^r,4\cdot6^r),
\]
with the same minimal period four.
\end{corollary}
Every base periodic orbit corresponds bijectively to a scaled periodic orbit of the same minimal period, and conversely.

Suppose $x_0,\ldots,x_{L-1}$ is a period-$L$ orbit in $D_r$. Put $a_i=1+e_r(x_i)$ and $c_i=c(x_i)$. Then
\[
5^{a_i}x_{i+1}=6x_i+6^rc_i.
\]
Let $A_j=\sum_{i=0}^{j-1}a_i$, $A_0=0$, and $A=A_L$.
\begin{theorem}[Scaled cycle identity]
Every period-$L$ orbit satisfies
\[
(5^A-6^L)x_0=6^r\sum_{j=0}^{L-1}6^{L-1-j}c_j5^{A_j}.
\]
Writing $x_0=6^rn_0$ reduces this identity exactly to the base cycle identity.
\end{theorem}
\begin{proof}Successive substitution gives the identity; division by $6^r$ gives the normalized form.\end{proof}
Since the right side is positive, every periodic orbit still requires $5^A>6^L$.

\section{Inverse branches and predecessor-tree transport}
Let $y\in D_r$ and suppose $T_r(x)=y$ with full denominator exponent $a\ge1$. Then
\[
5^ay=6x+6^rc,
\qquad x=\frac{5^ay-6^rc}{6},\quad c\in\{1,2,3,4\}.
\]
Writing $y=6^rz$ yields
\[
x=6^{r-1}(5^az-c).
\]
Membership $x\in D_r$ is equivalent to divisibility of $5^az-c$ by $6$, exactly the normalized base inverse-branch condition. Hence multiplication by $6^r$ gives a bijection between the complete predecessor tree of $z$ and that of $6^rz$ inside the scaled domains.

\section{Orbit statistics and normalized geometry}
Let $n_j=T_0^j(n)$ and $x_j=T_r^j(6^rn)=6^rn_j$.
Every order relation is preserved. Therefore first-descent time is invariant, first-lower values scale by $6^r$, elementary-step counts are invariant because the $5$-adic valuation word is unchanged, and peak positions are invariant. Moreover
\[
\max_{j\le J}x_j=6^r\max_{j\le J}n_j,
\qquad \frac{x_j}{x_0}=\frac{n_j}{n_0}.
\]
Thus absolute peaks scale by $6^r$, while peak/start ratios are identical.

\begin{theorem}[Homogeneous statistic transport]
If a finite-trajectory statistic $F$ satisfies $F(cz_0,\ldots,cz_J)=c^dF(z_0,\ldots,z_J)$, then corresponding scaled and base segments satisfy
\[
F(x_0,\ldots,x_J)=6^{rd}F(n_0,\ldots,n_J).
\]
\end{theorem}

\section{$5$-adic density and drift preservation}
The base map has the exact natural residue-density law
\[
\Pr(e=k)=\frac4{5^{k+1}},\qquad k\ge0,
\]
with $\mathbb E[e]=1/4$. On the sparse lattice $D_r$, the appropriate comparison is relative density after normalization $x=6^rn$. Since the valuation is preserved pointwise, every fixed finite valuation condition has exactly the same relative density at every scale.

For a matched step,
\[
\frac{T_r(x)}x=\frac{6^rT_0(n)}{6^rn}=\frac{T_0(n)}n.
\]
Hence normalized one-step multipliers and logarithmic increments are identical, not merely asymptotically similar. The base residue model therefore retains the same approximate multiplier $6/5^{e+1}$, arithmetic mean $1$, and negative mean logarithmic drift $\log6-(5/4)\log5$ at every scale after normalization.

\section{Transfer of the exhaustive verification through $10^{11}$}
For fixed $r$ define the corresponding scaled finite set
\[
S_r(10^{11})=\{6^rn:1\le n\le10^{11}\}.
\]
The base computation established that every base start in this interval enters $C$ and every $n\ge2$ has a strict first descent in at most $422$ macro steps.

\begin{theorem}[Finite verification transfer]\label{thm:finite}
For every fixed $r\ge0$, every state $x=6^rn$ with $1\le n\le10^{11}$ follows the base trajectory multiplied by $6^r$. In particular it enters $C_r$, and for $n\ge2$ its first-descent time is at most $422$ macro steps.
\end{theorem}
\begin{proof}Apply all-iterate conjugacy to each already verified base start. Strict inequalities are preserved by multiplication by the positive constant $6^r$.\end{proof}

\begin{remark}
Theorem~\ref{thm:finite} is not a claim that $10^{11}$ trajectories were independently recomputed for each $r$. It is an algebraic transfer of a completed finite computation through a proved conjugacy.
\end{remark}

The principal inherited records therefore become
\begin{center}
\begin{tabular}{lll}\toprule
Statistic & Invariant quantity & Scaled magnitude\\\midrule
First-descent record & $422$ steps & start $6^r(55,335,461,598)$\\
Elementary record & $470$ steps & start $6^r(55,335,461,598)$\\
Deepest observed valuation & $e=16$ & start $6^r(88,717,666,380)$\\
Largest pre-descent peak & peak position invariant & $6^r(8,669,800,813,122,806,007,604)$\\
Peak-producing start & --- & $6^r(37,266,730,278)$\\\bottomrule
\end{tabular}
\end{center}
The normalized peak/start ratio is unchanged.

\section{Explicit scaled trajectories}
The base cycle itself gives the simplest exact family:
\[
1\to2\to3\to4\to1
\]
becomes
\[
6^r\to2\cdot6^r\to3\cdot6^r\to4\cdot6^r\to6^r.
\]
For $r=1$ this is $6\to12\to18\to24\to6$; for $r=2$, $36\to72\to108\to144\to36$. These examples are direct consequences of the new definition and do not require independent search.

\section{Why the scaled domain is essential}
The expression $(6x+c(x)6^r)/5$ can be evaluated for many integers outside $D_r$, but the conjugacy theorem requires $x=6^rn$. Without this condition there is no factorization
\[
6x+6^rc(x)=6^r(6n+c(n)).
\]
Therefore the generalized affine formula on arbitrary integers should not be identified with a base orbit multiplied by $6^r$.

\begin{example}
Take $r=1$ and $x=1$, which is not in $D_1=6D$. Since $c(1)=4$,
\[
H_1(1)=\frac{6+24}{5}=6,
\]
whereas $6T_0(1)=12$. Thus the scaled-domain hypothesis is necessary.
\end{example}

\section{Global convergence and counterexample transport}
\begin{conjecture}[Base global convergence]
Every positive integer eventually enters $C$ under the base reduced map.
\end{conjecture}
\begin{conjecture}[Scaled global convergence]
For fixed $r$, every state in $D_r$ eventually enters $C_r$.
\end{conjecture}
\begin{theorem}[Equivalence across scales]
Base global convergence on $D$ holds if and only if scaled global convergence holds on $D_r$ for any fixed $r$. Nontrivial cycles and nonconvergent orbits likewise transport bijectively between scale levels.
\end{theorem}
This theorem preserves the unresolved status of the global problem: scaling neither proves nor disproves the base conjecture.

\section{Counting at corresponding cutoffs}
Let $A_0(X)=\#\{n<X:n\in D\}$ and $A_r(Y)=\#\{x<Y:x\in D_r\}$. The bijection gives
\[
A_r(6^rX)=A_0(X),\qquad A_r(Y)=A_0(Y/6^r).
\]
Thus exact transfer concerns corresponding normalized cutoffs, not equal absolute cutoffs.

\section{Computational consistency checks and reproducibility package}
The accompanying verifier checks, over a user-selected normalized range and scale levels,
\[
x\bmod5=n\bmod5,\quad H_r(6^rn)=6^rG_0(n),\quad e_r(6^rn)=e_0(n),\quad T_r(6^rn)=6^rT_0(n),
\]
and independently compares the residue-affine and floor-function forms. The package also contains machine-readable definitions, the base summary and selected base record files as provenance, and SHA-256 checksums. The verifier is a consistency aid; the conjugacy proof is algebraic and does not depend on it.

\section{What is proved, what is inherited, and what remains heuristic}
\paragraph{Exact algebra.} The generalized floor representation, residue preservation, affine factorization, $5$-adic valuation invariance, conjugacy, all-iterate correspondence, functional-graph isomorphism, scaled cycle identity, predecessor-tree transport, orbit-statistic transport, and equivalence of global questions are exact for all $r\ge0$ on the scaled domains.

\paragraph{Inherited finite theorem.} The completed base scan through $10^{11}$ transfers exactly to $S_r(10^{11})$ for every fixed $r$. Record times and valuation words are invariant; magnitude records acquire the factor $6^r$.

\paragraph{Probabilistic interpretation.} The one-step residue-density law is exact, but treating successive orbit states as independent or sufficiently mixing remains heuristic. The drift discussion is explanatory and is not a proof of universal convergence.

\section{Conclusion}
The family generated by
\[
\frac{6x+6^r}{5},\quad\frac{6x+2\cdot6^r}{5},\quad\frac{6x+3\cdot6^r}{5},\quad\frac{6x+4\cdot6^r}{5}
\]
according to the residue of $x$ modulo $5$, followed by complete $5$-adic reduction, has a simple exact interpretation on $D_r=6^rD$. The identity $6^r\equiv1\pmod5$ preserves the branch, while $\gcd(6^r,5)=1$ preserves the complete valuation. The result is the exact conjugacy $T_r(6^rn)=6^rT_0(n)$.

The requested floor form is given by Theorem~\ref{thm:floor}; at $r=0$ it collapses to $\lfloor6x/5\rfloor+1$. At every scale, the normalized dynamics are therefore identical to the base dynamics. The four-cycle scales to $(6^r,2\cdot6^r,3\cdot6^r,4\cdot6^r)$; cycles and predecessor trees correspond; stopping times and valuation words are unchanged; and absolute magnitudes scale by $6^r$.

The exhaustive base verification through $10^{11}$ consequently yields a simultaneous finite theorem on all corresponding scaled sets without repeating the massive computation. The global convergence question, however, remains exactly the base global problem in normalized coordinates.

\appendix
\section{Compact derivation of the scaling mechanism}
The entire argument is summarized by
\[
6^r\equiv1\pmod5\Rightarrow c(6^rn)=c(n),
\]
\[
6(6^rn)+6^rc(n)=6^r(6n+c(n)),
\]
\[
5\nmid6^r\Rightarrow\vfive(6^rm)=\vfive(m),
\]
\[
T_r(6^rn)=6^rT_0(n)\Rightarrow T_r^j(6^rn)=6^rT_0^j(n).
\]
Everything else in the scaled-domain theory follows from these identities and the positivity of the scale factor.

\section{Reproducibility commands}
A lightweight verification run is
\begin{verbatim}
python3 src/verify_scaled_family.py --limit 10000 --max-r 5 \
  --out data/verification_results.json
\end{verbatim}
The distributed JSON file records a passing run with 300,000 assertions.

\begin{thebibliography}{9}
\bibitem{Lagarias1985} J. C. Lagarias, ``The $3x+1$ problem and its generalizations,'' \emph{American Mathematical Monthly}, 92(1), 3--23, 1985. doi:10.1080/00029890.1985.11971528.
\bibitem{MatthewsWatts1985} K. Matthews and A. Watts, ``A Markov approach to the generalized Syracuse algorithm,'' \emph{Acta Arithmetica}, 45(1), 29--42, 1985. doi:10.4064/aa-45-1-29-42.
\bibitem{BanazadehBase} F. Banazadeh, ``The Reduced $\lfloor6k/5\rfloor+1$ /5 Collatz-Type Domain with One Observed Attractor up to $10^{11}$: Complete 5-Adic Reduction, Residue Algebra, Critical Multiplicative Statistics, Cycle Constraints, and Exhaustive First-Descent Verification,'' Zenodo, 2026, record 22856438.
\end{thebibliography}
\end{document}
